% Options for packages loaded elsewhere
% Options for packages loaded elsewhere
\PassOptionsToPackage{unicode}{hyperref}
\PassOptionsToPackage{hyphens}{url}
\PassOptionsToPackage{dvipsnames,svgnames,x11names}{xcolor}
%
\documentclass[
  12pt,
]{article}
\usepackage{xcolor}
\usepackage[left=2.5cm,right=2.5cm,top=2.5cm,bottom=2.5cm]{geometry}
\usepackage{amsmath,amssymb}
\setcounter{secnumdepth}{-\maxdimen} % remove section numbering
\usepackage{iftex}
\ifPDFTeX
  \usepackage[T1]{fontenc}
  \usepackage[utf8]{inputenc}
  \usepackage{textcomp} % provide euro and other symbols
\else % if luatex or xetex
  \usepackage{unicode-math} % this also loads fontspec
  \defaultfontfeatures{Scale=MatchLowercase}
  \defaultfontfeatures[\rmfamily]{Ligatures=TeX,Scale=1}
\fi
\usepackage{lmodern}
\ifPDFTeX\else
  % xetex/luatex font selection
\fi
% Use upquote if available, for straight quotes in verbatim environments
\IfFileExists{upquote.sty}{\usepackage{upquote}}{}
\IfFileExists{microtype.sty}{% use microtype if available
  \usepackage[]{microtype}
  \UseMicrotypeSet[protrusion]{basicmath} % disable protrusion for tt fonts
}{}
\makeatletter
\@ifundefined{KOMAClassName}{% if non-KOMA class
  \IfFileExists{parskip.sty}{%
    \usepackage{parskip}
  }{% else
    \setlength{\parindent}{0pt}
    \setlength{\parskip}{6pt plus 2pt minus 1pt}}
}{% if KOMA class
  \KOMAoptions{parskip=half}}
\makeatother
% Make \paragraph and \subparagraph free-standing
\makeatletter
\ifx\paragraph\undefined\else
  \let\oldparagraph\paragraph
  \renewcommand{\paragraph}{
    \@ifstar
      \xxxParagraphStar
      \xxxParagraphNoStar
  }
  \newcommand{\xxxParagraphStar}[1]{\oldparagraph*{#1}\mbox{}}
  \newcommand{\xxxParagraphNoStar}[1]{\oldparagraph{#1}\mbox{}}
\fi
\ifx\subparagraph\undefined\else
  \let\oldsubparagraph\subparagraph
  \renewcommand{\subparagraph}{
    \@ifstar
      \xxxSubParagraphStar
      \xxxSubParagraphNoStar
  }
  \newcommand{\xxxSubParagraphStar}[1]{\oldsubparagraph*{#1}\mbox{}}
  \newcommand{\xxxSubParagraphNoStar}[1]{\oldsubparagraph{#1}\mbox{}}
\fi
\makeatother
\usepackage{listings}
\newcommand{\passthrough}[1]{#1}
\lstset{defaultdialect=[5.3]Lua}
\lstset{defaultdialect=[x86masm]Assembler}


\usepackage{longtable,booktabs,array}
\usepackage{calc} % for calculating minipage widths
% Correct order of tables after \paragraph or \subparagraph
\usepackage{etoolbox}
\makeatletter
\patchcmd\longtable{\par}{\if@noskipsec\mbox{}\fi\par}{}{}
\makeatother
% Allow footnotes in longtable head/foot
\IfFileExists{footnotehyper.sty}{\usepackage{footnotehyper}}{\usepackage{footnote}}
\makesavenoteenv{longtable}
\usepackage{graphicx}
\makeatletter
\newsavebox\pandoc@box
\newcommand*\pandocbounded[1]{% scales image to fit in text height/width
  \sbox\pandoc@box{#1}%
  \Gscale@div\@tempa{\textheight}{\dimexpr\ht\pandoc@box+\dp\pandoc@box\relax}%
  \Gscale@div\@tempb{\linewidth}{\wd\pandoc@box}%
  \ifdim\@tempb\p@<\@tempa\p@\let\@tempa\@tempb\fi% select the smaller of both
  \ifdim\@tempa\p@<\p@\scalebox{\@tempa}{\usebox\pandoc@box}%
  \else\usebox{\pandoc@box}%
  \fi%
}
% Set default figure placement to htbp
\def\fps@figure{htbp}
\makeatother





\setlength{\emergencystretch}{3em} % prevent overfull lines

\providecommand{\tightlist}{%
  \setlength{\itemsep}{0pt}\setlength{\parskip}{0pt}}



 


\newcommand{\myCourse}{Bayesian Computational Methods}
\newcommand{\myAssignment}{Assignment }
\newcommand{\myTitle}{Help}
\newcommand{\myAuthor}{Isabella Lethbridge}
\newcommand{\myStudentNumber}{LTHISA001}
\newcommand{\myDate}{\today}
\usepackage{amsmath}
\usepackage{amssymb}
\usepackage{bm}
\usepackage{fancyhdr}
\usepackage{graphicx}
\usepackage{float}
\usepackage{subcaption}
\usepackage{etoolbox}

\usepackage{listings}
\usepackage{xcolor}

\renewcommand{\boldsymbol}[1]{\text{\pmb{$#1$}}}

% Shrink space around display math equations ($$ ... $$)
\AtBeginDocument{
  \setlength{\abovedisplayskip}{2pt}
  \setlength{\belowdisplayskip}{4pt}
  \setlength{\abovedisplayshortskip}{0pt}
  \setlength{\belowdisplayshortskip}{2pt}
}

\setlength{\abovedisplayskip}{4pt}
\setlength{\belowdisplayskip}{4pt}
\setlength{\abovedisplayshortskip}{2pt}
\setlength{\belowdisplayshortskip}{2pt}

% Define custom colors to match your target layout
\definecolor{codegreen}{rgb}{0,0.6,0}
\definecolor{codegray}{rgb}{0.5,0.5,0.5}
\definecolor{codepurple}{rgb}{0.58,0,0.82}
\definecolor{backcolour}{rgb}{0.95,0.95,0.92}
\definecolor{codeblue}{rgb}{1.0, 0.08, 0.58} % DeepPink

% Define R language rules for listings
\lstdefinelanguage{R}{
  keywords={TRUE,FALSE,NULL,NA,NaN,Inf,function,in,if,else,for,while,repeat,break,next},
  keywordstyle=\color{magenta}\bfseries,
  otherkeywords={=, <-, +,-,*,/,^\%,~\%,<,>,<=,>=,==,!=,!,&,|},
  morekeywords=[2]{c, data.frame, betareg, coef, vcov, cor, length, log, sum, pnorm, qnorm, nlm, library, require, sample, barplot, boxplot, prop.table, table, legend, colnames, t, diag, solve},
  keywordstyle=[2]\color{codeblue},
  sensitive=true,
  morecomment=[l]{\#},
  commentstyle=\color{codegreen},
  morestring=[b]",
  morestring=[b]',
  stringstyle=\color{codepurple}
}

% Apply the style globally
\lstdefinestyle{mystyle}{
  language=R,
  backgroundcolor=\color{backcolour},
  basicstyle=\ttfamily\footnotesize,
  breakatwhitespace=false,
  breaklines=true,
  captionpos=b,
  keepspaces=true,
  numbers=left,
  numbersep=5pt,
  numberstyle=\tiny\color{codegray},
  showspaces=false,
  showstringspaces=false,
  showtabs=false,
  tabsize=2
}

\lstset{style=mystyle}

% Append a page break directly to the \tableofcontents command
\gappto\tableofcontents{\clearpage}

\pagestyle{fancy}
\fancyhf{}
\lhead{\myTitle}
\cfoot{\thepage}
\makeatletter
\@ifpackageloaded{caption}{}{\usepackage{caption}}
\AtBeginDocument{%
\ifdefined\contentsname
  \renewcommand*\contentsname{Table of contents}
\else
  \newcommand\contentsname{Table of contents}
\fi
\ifdefined\listfigurename
  \renewcommand*\listfigurename{List of Figures}
\else
  \newcommand\listfigurename{List of Figures}
\fi
\ifdefined\listtablename
  \renewcommand*\listtablename{List of Tables}
\else
  \newcommand\listtablename{List of Tables}
\fi
\ifdefined\figurename
  \renewcommand*\figurename{Figure}
\else
  \newcommand\figurename{Figure}
\fi
\ifdefined\tablename
  \renewcommand*\tablename{Table}
\else
  \newcommand\tablename{Table}
\fi
}
\newcommand*\listoflistings\lstlistoflistings
\AtBeginDocument{%
\renewcommand*\lstlistlistingname{List of Listings}
}
\makeatother
\makeatletter
\makeatother
\makeatletter
\@ifpackageloaded{caption}{}{\usepackage{caption}}
\@ifpackageloaded{subcaption}{}{\usepackage{subcaption}}
\makeatother
\usepackage{bookmark}
\IfFileExists{xurl.sty}{\usepackage{xurl}}{} % add URL line breaks if available
\urlstyle{same}
\hypersetup{
  colorlinks=true,
  linkcolor={blue},
  filecolor={Maroon},
  citecolor={Blue},
  urlcolor={blue},
  pdfcreator={LaTeX via pandoc}}


\author{}
\date{}
\begin{document}

\begin{titlepage}
  \centering
  \vspace*{1.5cm}
  
  \includegraphics[width=0.35\textwidth]{UCT.jpg}\\[1.5cm]
  
  {\Large \textsc{University of Cape Town}}\\[2.0cm]
  
  {\large \textsc{\myCourse}}\\[0.3cm]
  {\large \textsc{\myAssignment}}\\[0.8cm]
  
  \rule{\linewidth}{0.3mm}\\[0.5cm]
  {\huge \bfseries \myTitle}\\[0.4cm]
  \rule{\linewidth}{0.3mm}\\[1.6cm]
  
  \begin{minipage}{0.45\textwidth}
    \begin{flushleft} \large
      \emph{Author:}\\
      \myAuthor
    \end{flushleft}
  \end{minipage}
  \begin{minipage}{0.45\textwidth}
    \begin{flushright} \large
      \emph{Student Number:}\\
      \myStudentNumber
    \end{flushright}
  \end{minipage}\\[3.0cm]
  
  {\large \myDate}\\[1cm]
  \vfill
\end{titlepage}
\clearpage


\subsection{Introduction}\label{introduction}

\subsection{Question a}\label{question-a}

We can write the probability that \(s_i=3\) as

\[
P(t_i) = Pr(T_i > t_i) = \int^{\infty}_{t_i} \frac{1}{\sqrt{2 \pi} \sigma} \text{exp} \big( -\frac{1}{2\sigma^2} (x - \mu)^2\big) dx.
\]

This can be written in terms of \(\Phi\) after standardising

\[
P(t_i) = Pr(\frac{T_i-\mu}{\sigma} > \frac{t_i-\mu}{\sigma}) = 1 - \Phi(\frac{t_i-\mu}{\sigma}).
\]

The probability that \(s_i=2\) is denoted as

\[
Q(t_i) = Pr(t_i - \tau \leq T_i \leq t_i) = \int^{t_i}_{t_i - \tau} \frac{1}{\sqrt{2 \pi} \sigma} \text{exp} \big( -\frac{1}{2\sigma^2} (x - \mu)^2\big) dx.
\]

This can be written in terms of \(\Phi\) after standardising

\[
Q(t_i) = Pr(\frac{t_i - \tau - \mu}{\sigma} \leq\frac{T_i-\mu}{\sigma} \leq \frac{t_i-\mu}{\sigma}) = \Phi(\frac{t_i-\mu}{\sigma}) - \Phi(\frac{t_i -\tau - \mu}{\sigma}).
\]

The probability that \(s_i=1\) is denoted as

\[
R(t_i) = Pr(T_i \leq t_i - \tau) = \int^{t_i - \tau}_{- \infty} \frac{1}{\sqrt{2 \pi} \sigma} \text{exp} \big( -\frac{1}{2\sigma^2} (x - \mu)^2\big) dx.
\]

This can be written in terms of \(\Phi\) after standardising

\[
R(t_i) = Pr(\frac{T_i-\mu}{\sigma} \leq \frac{t_i-\tau-\mu}{\sigma}) = \Phi(\frac{t_i-\tau-\mu}{\sigma}).
\]

We now have

\[
L(\boldsymbol{s}, \boldsymbol{t}, \mu, \sigma^2, \tau) = \prod^{n}_{i=1} [1 - \Phi(\frac{t_i-\mu}{\sigma})]^{\text{I}(s_i=3)} \cdot [\Phi(\frac{t_i-\mu}{\sigma}) - \Phi(\frac{t_i -\tau - \mu}{\sigma})]^{\text{I}(s_i=2)} \cdot [\Phi(\frac{t_i-\tau-\mu}{\sigma})]^{\text{I}(s_i=1)}.
\]

\subsection{Question b}\label{question-b}

Conditional posterior distributions to iteratively sample from using a
Gibbs sampler:

\[
\begin{aligned}
T_i \mid \mu,\sigma^2,\tau,s_i =1, \quad t_i & \sim \text{N}(\mu, \sigma^2) I_{\{T_i \leq t_i - \tau \}} \forall i \\
T_i \mid \mu,\sigma^2,\tau,s_i =2, \quad t_i & \sim \text{N}(\mu, \sigma^2) I_{\{t_i-\tau \leq T_i \leq t_i \}} \forall i \\
T_i \mid \mu,\sigma^2,\tau,s_i =3, \quad t_i & \sim \text{N}(\mu, \sigma^2) I_{\{T_i > t_i \}} \forall i \\
\mu \mid \sigma^2,\tau, \boldsymbol{s} , \boldsymbol{t} , \boldsymbol{T} & \sim \text{N} \big(\bar{T}, \frac{\sigma^2}{n} \big) I_{\{100 \leq \mu \leq 160 \}} \\
\sigma^2 \mid \mu, \tau, \boldsymbol{s}, \boldsymbol{t}, \boldsymbol{T} & \sim \text{IG} \big( \frac{n}{2}, \frac{1}{2} \sum^n_{i=1} (T_i - \mu)^2 \big) \\
\tau \mid . & \sim \text{Unif}(\tau \in V)
\end{aligned}
\]

where \(\text{V} = [50,180] \cap [l,h]\)

Under the Type 2 model, one has observed the moult scores for the birds
undergoing moult. Fot those birds the probability density function of
the moult score are

\[
f_y(y) = \tau f_T(t-y\tau) \qquad 0<y<1. 
\]

After observing the moult scores when \(s_i=2\) we need to update the
conditional distributions. Only
\(T_i \mid \mu,\sigma^2,\tau,s_i =2, \quad t_i  \sim \text{N}(\mu, \sigma^2) I_{\{t_i-\tau \leq T_i \leq t_i \}} \forall i\)
and \(\tau \mid . \sim \text{Unif}(\tau \in V)\) change.

We know that \(T_i = t_i-y_i \cdot \tau\).

We can write the probability density function (ref eq here) as (for a
single obervation)

\[
\begin{aligned}
f_y(y) & = \tau \cdot \frac{1}{\sqrt{2\pi}\sigma} \text{exp}\big(-\frac{1}{2\sigma^2}(t-y\cdot \tau-\mu)^2\big) \\
& = \tau \cdot \frac{1}{\sqrt{2\pi}\sigma} \text{exp}\big(-\frac{1}{2\sigma^2}(t^2 - 2ty \tau + (y\tau)^2 - 2\mu t + 2\mu y \tau +\mu^2\big) \\
& \propto \tau \cdot \text{exp}\big(-\frac{1}{2\sigma^2}((y\tau)^2 -2y\tau(t-\mu))\big) \\
& \propto \tau \cdot \text{exp} \big( -\frac{y^2}{2\sigma^2} \tau^2 + \frac{y(t-\mu)}{\sigma^2}\tau\big)
\end{aligned}
\]

The writing it as a likelihood (remember \(t_i\) changes for each \(i\))

\[
\begin{aligned}
& \prod^m_{i=1} \tau \cdot \text{exp} \big( -\frac{y_i^2}{2\sigma^2} \tau^2 + \frac{y_i(t_i-\mu)}{\sigma^2}\tau\big) \\
= & \tau^m \cdot \text{exp}\big(-(\frac{1}{2\sigma^2} \sum^m_{i=1} y_i^2)\tau^2 + (\frac{1}{\sigma^2} \sum^m_{i=1} y_i(t_i-\mu))\tau\big)
\end{aligned}
\]

According to Lemma 5 of Sun et al.~(2021)

\[
\text{Let } X \sim \text{MHN}(\alpha, \beta, \gamma)
\]

\begin{enumerate}
\def\labelenumi{\alph{enumi}.}
\item
  If \(\gamma>0\) then there exists a random variable \(V\) such that
  \(V\mid X \sim \text{Poisson}(\gamma X)\) and
  \(X^2 \mid V \sim \text{Gamma}(\frac{\alpha +V}{2}, \beta)\).
\item
  If \(\gamma<0\) then there exists a random variable \(U\) such that
  \(U\mid X \sim \text{GIG}(\frac{1}{2},1,\gamma^2 X^2)\) and
  \(X^2 \mid U \sim \text{Gamma}(\frac{\alpha}{2}, (\beta + \frac{\gamma^2}{U})\).
\end{enumerate}

The PDF for the modified half normal is

\[
\begin{aligned}
f(x\mid \alpha, \beta, \gamma) & = \dfrac{2 \beta^{\alpha/2} x^{\alpha-1} \text{exp}(-\beta x^2 + \gamma x)}{\Psi [\frac{\alpha}{2}, \frac{\gamma}{\sqrt{\beta}}]} \\
& \propto x^{\alpha-1} \text{exp}(-\beta x^2 + \gamma x) 
\end{aligned}
\]

From (ref eq here) we have \(\alpha = m+1\),
\(\beta = \frac{1}{2\sigma^2} \sum^m_{i=1} y_i^2\) and
\(\gamma = \frac{1}{\sigma^2} \sum^m_{i=1} y_i(t_i-\mu)\).

We can now write the distribution for \(\tau\) as (Notation is getting a
bit confusing here as we now have 2 V's)

\[
\tau \sim \text{MHN}\big(\alpha = m+1, \beta = \frac{1}{2\sigma^2} \sum^m_{i=1} y_i^2, \gamma = \frac{1}{\sigma^2} \sum^m_{i=1} y_i(t_i-\mu)\big) \cdot I_{\{\tau \in \text{V}\}}.
\]

From \(s_i=1\) we have

\[
\begin{aligned}
T_i & \leq t_i - \tau \\
\tau & \leq t_i - T_i
\end{aligned}
\]

Since \(t_i\) changes for every \(i\) we have the upper bound
\(h = \text{min}(t_i -T_i) \forall i\). (I'm not sure if this is the
most elegant way of writing this but we can fix it later).

For \(s_i = 3\) we have

\[
\begin{aligned}
t_i - \tau & \leq T_i \leq t_i \\
\tau & \geq t_i - T_i \geq 0 \\
\tau & \geq t_i - (t_i - y_i \tau) \\
\tau & \geq y_i\tau
\end{aligned}
\]

Since \(0<y_i<1\) and \(\tau>0\), this is always true.

Thus \(\text{V} = [50,min(180,h)]\) and
\(\tau \in [50^2, \text{min}(100,h)^2]\).

We now consider case (a): \(\gamma > 0\)

Let \(m\) be the number of birds with \(s_i=2\) (the moutling birds).

\[
\begin{aligned}
\gamma = \frac{1}{\sigma^2} \sum^m_{i=1} y_i(t_i-\mu) & > 0 \\
\sum^m_{i=1} y_i(t_i-\mu) & > 0
\end{aligned}
\]

\[
\begin{aligned}
V \mid \tau & \sim \text{Poisson}(\gamma \cdot \tau) \\
& = \text{Poisson}\big(\frac{\tau}{\sigma^2} \sum^m_{i=1} y_i(t_i-\mu)\big)
\end{aligned}
\]

\[
\tau^2 \mid V \sim \text{Gamma}\big(\frac{m+1+V}{2},\frac{1}{2\sigma^2} \sum^m_{i=1} y_i^2\big)
\]

Next consider case (b): \(\gamma < 0\)

\[
\begin{aligned}
\gamma = \frac{1}{\sigma^2} \sum^m_{i=1} y_i(t_i-\mu) & < 0 \\
\sum^m_{i=1} y_i(t_i-\mu) & < 0
\end{aligned}
\]

\[
U \mid \tau \sim \text{GIG}\big(\frac{1}{2},1,(\frac{\tau}{\sigma^2} \sum^m_{i=1} y_i(t_i-\mu))^2\big)
\]

\[
\tau^2 \mid U \sim \text{Gamma}\big(\frac{m+1}{2}, (\frac{1}{2\sigma^2} \sum^m_{i=1} y_i^2 + \frac{1}{U}\cdot(\frac{1}{\sigma^2} \sum^m_{i=1} y_i(t_i-\mu))^2)\big)
\]

Now that we have all these lovely distributions we actually need to do
this data augmentation shit

\textbf{We must check the parameterisation of the Gamma dist}

\subsubsection{Gibbs sampler for the Type 2
model}\label{gibbs-sampler-for-the-type-2-model}

Let \(m\) be the number of birds undergoing moult (\(s_i = 2\)). The
sets \(\{i: s_i = 1\}\), \(\{i: s_i = 2\}\) and \(\{i: s_i = 3\}\)
denote the birds that have completed moult, are undergoing moult and
have not yet started moult respectively. For the birds undergoing moult
the start of moult is not latent anymore since it is fully determined by
\(\tau\) and the observed moult score,

\[
T_i = t_i - y_i \tau, \qquad i: s_i = 2.
\]

The latent start dates therefore only need to be sampled for the birds
with \(s_i = 1\) and \(s_i = 3\). The conditional distribution of
\(\tau\) is no longer uniform since \(\tau\) now appears in the density
of the observed moult scores. From the derivation above, the conditional
of \(\tau\) is a truncated MHN distribution,

\[
\tau \mid \mu, \sigma^2, \boldsymbol{T}, \boldsymbol{y}, \boldsymbol{s}, \boldsymbol{t} \sim \text{MHN}(\alpha, \beta, \gamma) \, I_{\{\tau_L \leq \tau \leq \tau_U\}}
\]

with

\[
\alpha = m + 1, \qquad \beta = \frac{1}{2\sigma^2} \sum_{i: s_i = 2} y_i^2, \qquad \gamma = \frac{1}{\sigma^2} \sum_{i: s_i = 2} y_i (t_i - \mu).
\]

The uniform prior on \(\tau\) and the birds that have completed moult
restrict \(\tau\) to the interval

\[
\tau_L = 50, \qquad \tau_U = \min\Big(180, \min_{i: s_i = 1} (t_i - T_i)\Big).
\]

The birds that have not started moult do not restrict \(\tau\) since
\(\tau\) does not appear in the condition \(T_i > t_i\). The birds
undergoing moult also do not restrict \(\tau\) anymore, since
substituting \(T_i = t_i - y_i\tau\) into
\(t_i - \tau \leq T_i \leq t_i\) gives \(y_i \tau \leq \tau\) and
\(y_i \tau \geq 0\), which always holds for \(0 < y_i < 1\). This is
different from the Type 1 model, where the birds undergoing moult gave
the lower bound \(l\).

Sampling from the MHN directly is difficult, so we use data augmentation
(Lemma 5 of Sun et al.~(2021)). We introduce an auxiliary variable
(\(V\) if \(\gamma > 0\) or \(U\) if \(\gamma < 0\)) and sample it
together with \(\tau\). The marginal distribution of \(\tau\) is still
the MHN, but each of the two conditionals is a standard distribution
that is easy to sample from.

\textbf{Algorithm}

\emph{Initialisation.} For each chain draw random starting values from
the priors, \(\mu^{(0)} \sim \text{Uniform}(100, 160)\),
\(\sigma^{2(0)} \sim \text{InverseGamma}(1,1)\) and
\(\tau^{(0)} \sim \text{Uniform}(50, 180)\). Then set
\(T_i^{(0)} = t_i - y_i \tau^{(0)}\) for all \(i\) with \(s_i = 2\).

\emph{For iteration} \(k = 1, \dots, N\):

\begin{enumerate}
\def\labelenumi{\arabic{enumi}.}
\tightlist
\item
  \textbf{Birds that have completed moult.} For each \(i\) with
  \(s_i = 1\) sample
\end{enumerate}

\[
T_i^{(k)} \sim \text{N}\big(\mu^{(k-1)}, \sigma^{2(k-1)}\big) \, I_{\{T_i \leq t_i - \tau^{(k-1)}\}}.
\]

\begin{enumerate}
\def\labelenumi{\arabic{enumi}.}
\setcounter{enumi}{1}
\tightlist
\item
  \textbf{Birds that have not started moult.} For each \(i\) with
  \(s_i = 3\) sample
\end{enumerate}

\[
T_i^{(k)} \sim \text{N}\big(\mu^{(k-1)}, \sigma^{2(k-1)}\big) \, I_{\{T_i > t_i\}}.
\]

The truncated normal draws in steps 1 and 2 are done using the
\passthrough{\lstinline!RcppTN!} package.

\begin{enumerate}
\def\labelenumi{\arabic{enumi}.}
\setcounter{enumi}{2}
\tightlist
\item
  \textbf{Update} \(\mu\). Sample
\end{enumerate}

\[
\mu^{(k)} \sim \text{N}\Big(\bar{T}, \frac{\sigma^{2(k-1)}}{n}\Big) \, I_{\{100 \leq \mu \leq 160\}},
\]

where \(\bar{T}\) is the mean over all \(n\) birds, using the sampled
\(T_i^{(k)}\) for \(s_i = 1, 3\) and the calculated
\(T_i = t_i - y_i \tau^{(k-1)}\) for \(s_i = 2\).

\begin{enumerate}
\def\labelenumi{\arabic{enumi}.}
\setcounter{enumi}{3}
\tightlist
\item
  \textbf{Update} \(\sigma^2\). Sample
\end{enumerate}

\[
\sigma^{2(k)} \sim \text{IG}\Big(\frac{n}{2}, \frac{1}{2} \sum_{i=1}^n \big(T_i - \mu^{(k)}\big)^2\Big),
\]

again using all \(n\) start dates.

\begin{enumerate}
\def\labelenumi{\arabic{enumi}.}
\setcounter{enumi}{4}
\item
  \textbf{Update} \(\tau\) (data augmentation).

  \begin{enumerate}
  \def\labelenumii{\alph{enumii}.}
  \item
    Calculate \(\alpha = m+1\),
    \(\beta = \frac{1}{2\sigma^{2(k)}} \sum_{i: s_i = 2} y_i^2\) and
    \(\gamma = \frac{1}{\sigma^{2(k)}} \sum_{i: s_i = 2} y_i (t_i - \mu^{(k)})\).
  \item
    Calculate the bounds \(\tau_L = 50\) and
    \(\tau_U = \min\big(180, \min_{i: s_i = 1} (t_i - T_i^{(k)})\big)\).
  \item
    If \(\gamma > 0\), sample the auxiliary variable and then
    \(\tau^2\):
  \end{enumerate}
\end{enumerate}

\[
V \mid \tau^{(k-1)} \sim \text{Poisson}\big(\gamma \, \tau^{(k-1)}\big), \qquad \tau^2 \mid V \sim \text{Gamma}\Big(\frac{\alpha + V}{2}, \beta\Big) \, I_{\{\tau_L^2 \leq \tau^2 \leq \tau_U^2\}}.
\]

If \(\gamma < 0\), sample

\[
U \mid \tau^{(k-1)} \sim \text{GIG}\big(\tfrac{1}{2}, 1, \gamma^2 (\tau^{(k-1)})^2\big), \qquad \tau^2 \mid U \sim \text{Gamma}\Big(\frac{\alpha}{2}, \beta + \frac{\gamma^2}{U}\Big) \, I_{\{\tau_L^2 \leq \tau^2 \leq \tau_U^2\}}.
\]

\begin{enumerate}
\def\labelenumi{\alph{enumi}.}
\setcounter{enumi}{3}
\item
  Set \(\tau^{(k)} = \sqrt{\tau^2}\).
\item
  Recalculate \(T_i = t_i - y_i \tau^{(k)}\) for all \(i\) with
  \(s_i = 2\), so that the next iteration uses start dates that agree
  with the new value of \(\tau\).
\end{enumerate}

\begin{enumerate}
\def\labelenumi{\arabic{enumi}.}
\setcounter{enumi}{5}
\tightlist
\item
  \textbf{Store} \(\mu^{(k)}\), \(\sigma^{2(k)}\) and \(\tau^{(k)}\).
\end{enumerate}

The truncated Gamma in step 5c is sampled using the inverse CDF method.
Let \(F\) be the CDF of the (untruncated) Gamma distribution. We draw
\(u \sim \text{Uniform}\big(F(\tau_L^2), F(\tau_U^2)\big)\) and set
\(\tau^2 = F^{-1}(u)\). Every draw therefore lies in
\([\tau_L^2, \tau_U^2]\) and no rejection step is needed. The interval
is never empty, since the current \(\tau\) is at least 50 and the latent
\(T_i\) for \(s_i = 1\) were drawn so that \(\tau \leq t_i - T_i\).

We run three chains of 50 000 iterations each, every chain with its own
random starting values.

\subsection{Question c}\label{question-c}

We sample from the likelihood surface using a Metropolis Hastings
algorithm. However we have the problem that the steps of our random walk
are drawn a multivariate normal distribution. But each of our three
parameters do not have support over the entire real line: \(\mu\), which
determines the average start time, is constrained by a uniform prior to
be situated between 100--160; \(\tau\), which determines the speed of
moulting, is also contrainted by its prior to be situated within
50--180; finally the variance is strictly positive. We define the
following transformations:\\
\[\begin{align}
f(\mu)&=\text{logit} \frac{\mu-100}{60}\\
g(\tau)&=\text{logit} \frac{\tau-50}{130}\\
h(\sigma^2)&= \log \sigma^2
\end{align}\] These transformation each have support over the entire
real line; thus we may perform our random walk over this transformed
space. Since we are unsure how our parameters are correlated; to begin
with, the steps of our random walk are drawn from:\\
\[e_i\sim\mathcal{N}(0,c \,I)\] For some constant \(c\). However, when
back transform (let \(\theta\) be our original parameters and \(\phi\)
their transformations) we must include the jacobian term in the
posterior likelihood in order to ensure that the Metropolis-Hasting
algorithm converges to the appropriate likelihood; the log-likelihood
becomes: \[L(\theta)+\log(\pi(\theta))+\log|J(\phi)|\] Where \(J(\phi)\)
is the Jacobian of our transformation. However, we would like to tune
our random walk to have a more appropriate covaraince matrix; to achieve
this we utilize 3000 burnin iterations and estimate the covariance
matrix \(\Sigma\) of our transformed parameters on these iterations.
Them each of our retained draws each have their random walk drawn from:
\[e\sim\mathcal{N}(0,\Sigma)\] And the Metropolis Hastings algorithm
algorithm proceeds as per usual. We obtain the following estimates of
our parameters and as well as 95\% credible intervals on them:

\begin{longtable}[]{@{}
  >{\raggedright\arraybackslash}p{(\linewidth - 6\tabcolsep) * \real{0.3733}}
  >{\raggedleft\arraybackslash}p{(\linewidth - 6\tabcolsep) * \real{0.2133}}
  >{\raggedleft\arraybackslash}p{(\linewidth - 6\tabcolsep) * \real{0.2000}}
  >{\raggedleft\arraybackslash}p{(\linewidth - 6\tabcolsep) * \real{0.2133}}@{}}
\caption{Posterior mean estimates and 95\% credible intervals for the
Type 3 moulting model.}\label{tbl-mh-estimates}\tabularnewline
\toprule\noalign{}
\begin{minipage}[b]{\linewidth}\raggedright
Parameter
\end{minipage} & \begin{minipage}[b]{\linewidth}\raggedleft
Posterior mean
\end{minipage} & \begin{minipage}[b]{\linewidth}\raggedleft
2.5\% quantile
\end{minipage} & \begin{minipage}[b]{\linewidth}\raggedleft
97.5\% quantile
\end{minipage} \\
\midrule\noalign{}
\endfirsthead
\toprule\noalign{}
\begin{minipage}[b]{\linewidth}\raggedright
Parameter
\end{minipage} & \begin{minipage}[b]{\linewidth}\raggedleft
Posterior mean
\end{minipage} & \begin{minipage}[b]{\linewidth}\raggedleft
2.5\% quantile
\end{minipage} & \begin{minipage}[b]{\linewidth}\raggedleft
97.5\% quantile
\end{minipage} \\
\midrule\noalign{}
\endhead
\bottomrule\noalign{}
\endlastfoot
Mean start time (\(\mu\)) & 130.67 & 114.61 & 152.40 \\
Variance (\(\sigma\^2\)) & 317.18 & 195.19 & 525.00 \\
Moulting duration (\(\tau\)) & 92.79 & 56.43 & 128.40 \\
\end{longtable}




\end{document}
